\uuid{GJ8f}
\titre{Équation droite/plan}
\auteur{bodin ; hanani ; mzari}
\organisation{exo7}
\licence{CC BY-NC-SA 4.0}
\datecreate{2020-01-20}
\niveau{L1}
\module{Géométrie}
\chapitre{Géométrie du plan}
\sousChapitre{Distance, norme, produit scalaire}
\difficulte{2}
\theme{équation droite, plan}
\exo7qcmid{lille-1.geometrie-plane.017}

\contenu{
\begin{question}
\qclassification{140.01,140.02}
\qtags{motcle=équation droite/plan}
Dans le plan rapporté à un repère orthonormé, on considère la droite $D$ passant par les points $A(1,1)$ et $B(2,3)$. Quelles sont les bonnes réponses ?
\begin{answers}
    \bad{$\vec{u}=(1,2)$ est un vecteur normal à $D$.}
    \good{Une équation cartésienne de $D$ est : $2x-y-1=0$.}
    \bad{Le point $C(1,2)$ appartient à $D$.}
    \good{La distance du point $N(-1,2)$ à la droite $D$ est $\sqrt{5}$.}
\end{answers}
\begin{explanations}
Le vecteur $\overrightarrow{AB}=(1,2)$ dirige $D$ et $M(x,y)\in D\Leftrightarrow \mbox{det} \left(\overrightarrow{AM},\overrightarrow{AB}\right)=0$, c'est-à-dire $2x-y-1=0$. La distance de $N$ à $D$ est donnée par $\displaystyle \frac{|2\times (-1)-2-1|}{\sqrt{2^2+1^2}}=\sqrt{5}$.
\end{explanations}
\end{question}
}
